\chapter{Answers to ``Check yourself''}
\label{app:answers}

\begin{enumerate}[leftmargin=1.1cm]
\answitem{\textbf{(Ch.~\ref{ch:building})} An answer can be right by luck or on easy items
  while the stated confidence does not track correctness across the distribution. You catch it
  by measuring calibration (a reliability diagram and ECE against the noise floor), not just
  accuracy.}
\answitem{\textbf{(Ch.~\ref{ch:roadmap})} The record/replay cache makes every result
  reproducible offline with no API calls, so at gate review you can regenerate every table and
  audit exactly what was measured. It also removes network flakiness and cost from the
  build-test loop.}
\answitem{\textbf{(Ch.~\ref{ch:agents})} Guaranteed: valid output shape and a valid schema for
  tool arguments. Not guaranteed: a probability distribution over which tool was chosen. The
  decision layer adds that calibrated probability.}
\answitem{\textbf{(Ch.~\ref{ch:jev})} (1) Adversarial content is a documented Jev weakness, so
  it can be fooled by a benign-looking wrapper. (2) A saturated confidence of 0.99 can be
  wrong, and safety-critical authorization must be deterministic; code and allow/deny lists,
  not a probability, decide.}
\answitem{\textbf{(Ch.~\ref{ch:calibration})} 92\% accuracy with ECE 0.15 means the
  probabilities are badly miscalibrated, so ``confidence $\ge 0.9$'' may include many wrong
  items. Inspect the reliability diagram (Figure~\ref{fig:reliability}) and the ECE-vs-floor
  bars (Figure~\ref{fig:ece}) first, then recalibrate before gating.}
\answitem{\textbf{(Ch.~\ref{ch:routers})} Because Jev is so cheap it can be queried
  \emph{before} paying for the LLM, like a pre-generation router; a cascade instead pays the
  cheap model first and only then decides, which wastes the cheap call when it defers.}
\answitem{\textbf{(Ch.~\ref{ch:eval})} pass@8 counts a single lucky success out of eight,
  while pass\^{}8 requires success in all eight; an unreliable per-attempt rate (say 0.7)
  gives $0.7^8 \approx 0.058$. See Figure~\ref{fig:passk}.}
\answitem{\textbf{(Ch.~\ref{ch:decisionpoints})} DP2 narrows because Jev cannot write tool
  arguments (no generation) and its choice is advisory; keeping the LLM in the loop preserves
  correctness. DP4 never declares success alone because ``done'' is safety-relevant and must
  be confirmed by a deterministic completion verifier.}
\answitem{\textbf{(Ch.~\ref{ch:arch})} The call is treated as destructive (most conservative
  wins); DP3 output can only raise the required approval level, never lower it. The principle
  is ``the engine advises, code authorizes.''}
\answitem{\textbf{(Ch.~\ref{ch:engineering})} If the fallback reused the lossy projection while
  the baseline saw full context, fallback-heavy runs would look worse for reasons unrelated to
  the engine, biasing the whole cost/quality comparison. Rebuilding full context keeps the
  arms comparable.}
\answitem{\textbf{(Ch.~\ref{ch:evidence})} H5: if ARM-A2 (an LLM answering the same typed
  questions) matches ARM-A1, the benefit is from typed decision points, not Jev specifically,
  so the product story shifts from ``Jev'' to ``calibrated typed decisions, with Jev as the
  cheap option.''}
\end{enumerate}
